% PIC-tutor native LaTeX book — hyperlinks, labels, bibliography configuration,
% and the semantic reader-element macros (Phase 1 set).

% --- Colors (reused from the Pandoc-era theme palette for continuity) ---
\definecolor{PicTutorInk}{HTML}{202A33}
\definecolor{PicTutorAccent}{HTML}{355C7D}
\definecolor{PicTutorLink}{HTML}{1D5E86}
\definecolor{PicTutorCodeBorder}{HTML}{B8C4CE}

% --- hyperref / cleveref ---
\hypersetup{
  colorlinks=true,
  linkcolor=PicTutorAccent,
  urlcolor=PicTutorLink,
  citecolor=PicTutorLink,
  bookmarksnumbered=true,
  bookmarksopen=true,
  pdftitle={PIC-tutor},
}
\crefname{chapter}{第}{章}
\crefname{section}{第}{节}
\crefname{figure}{图}{}
\crefname{table}{表}{}
\crefname{equation}{式}{}

% --- Semantic reader-element macros (Phase 1 minimal set) ---
% Source-path label lines used throughout the book, e.g.
%   源码入口：`Source/WarpX.cpp`
%   源码文件：`Source/Particles/Pusher/UpdateMomentumBoris.H`
%   函数：`UpdateMomentumBoris()`
\newcommand{\sourceline}[2]{\par\noindent\textbf{#1：}#2\par}

% Standalone bold card labels, e.g. "**本章的读者路线**", "**读者的时间层判断卡**".
\newcommand{\cardlabel}[1]{\par\noindent\textbf{#1}\par}

% Card step paragraphs inside validation cards, e.g. "**第一层：...**".
\newcommand{\cardstep}[1]{\par\noindent\textbf{#1}\par}

% Internal chapter reference by canonical label (see converter label scheme).
% Falls back to plain text while the target chapter is not yet in the book.
\newcommand{\chref}[2]{\hyperref[#1]{#2}}
